\section{Reintroduction}\label{sec:reintroduction}

The journey from the psychoanalyst's couch to the neuroscience lab has fundamentally reshaped our understanding of sleep. It has clarified that ``dreaming'' is not a single state, but a collection of distinct phenomenological experiences (like hypnagogia and REM mentation) tied to specific neurobiological architectures \citep{revonsuo2013important}.

We now understand sleep as an active, multi-stage process critical for memory consolidation \citep{wilson1994reactivation} and cognitive maintenance. The severe cognitive and psychological consequences of sleep deprivation \citep{waters2018deprivation} are not merely a result of fatigue, but a catastrophic failure of this essential offline processing. This paper will argue that these established pillars of sleep science—its complex architecture, its role in memory, and its non-negotiable nature—can be best understood through a single, unifying computational framework.